\scanpage{26}
5. Modules libres de rang $2$ sur $\Lambda=\mathbb{Z}/4\mathbb{Z}$.

Soit $M$ un tel module, et soit
\[
(1)\qquad M_0=M/2M=M\otimes_\Lambda\mathbb{F}_2
\]
qui est un vectoriel de rang $2$ sur $\mathbb{F}_2$. Soit $J=M_0^*$, on utilise la notation indicielle pour les éléments de $M_0^*$, notés $x_i$ ($i\in J$). \deleted{Pour tout $i\in J$, soit}

On a une suite exacte de \deleted{groupes \unsure{comm.}}\inserted{$\Lambda$-modules}.
\diagram{\begin{tikzcd}
(2)\quad0 \arrow[r,dash] & M_0 \arrow[r] & M \arrow[r,"\mathrm{can}"] & M_0 \arrow[r] & 0\\
& M \arrow[u,dotted,"\mathrm{can}"'] \arrow[ur,dotted,"2"'] &&&
\end{tikzcd}}
donc $M$ apparaît comme extension de $M_0$ par $M_0$ (en tant que $\Lambda$-module). Pour $\forall x\in M_0$, soit $M_x$ son image inverse dans $M$, qui est un torseur sous $M_0$ (i.e. un ens. à quatre éléments \inserted{$M_x$}, muni d'un iso. $M_0\simeq V_{M_x}$).

On posera, pour $i\in J$
\[
(3)\qquad
\left\{\begin{array}{l}
M_i=M_{x_i}=\deleted{\text{\unsure{image inverse de $\{x_i\}$}}}
\inserted{\text{fibre de $M$ sur }}x_i\in M_0,\ M_0\text{-torseur}\\
M_J=\prod_{i\in J}M_i,\ \text{un torseur sous }\deleted{\unsure{\mathbb{F}_2}}M_0^J.
\end{array}\right.
\]
Comme on
\[
(4)\qquad\sum_{i\in J}x_i=0
\]
la structure de groupe sur $M$ donne un iso. can.
\[
(6)\qquad\bigwedge_{i\in J}M_i=0
\]
en d'autres termes, introduisant
\[
(7)\qquad V_J(M_0)=\left\{(\unsure{\xi_i|\bar{x}_i})_{i\in J}\in M_0^J\mid\sum\xi_i=0\right\}
\]
s'insérant dans la suite exacte
\diagram{\begin{tikzcd}
(8)\quad0 \arrow[r,dash] & V_J(M_0) \arrow[r,dash] & M_0^J \arrow[r,dash,"\text{somme}"'] & M_0 \arrow[r,dash] & 0,
\end{tikzcd}}
le torseur $M_J$ sur $M_0^J$ provient canoniquement d'un torseur sous $V_J(M_0)$, que je note $P_M$,
% archive foliation: 25
